% !PRINTED 419
\weberpage{419}
%<<CONT
\weq{3}{\varPhi(x) = \varphi(x)\,\varphi(x^2)\,\varphi(x^3)\dots\varphi(x^{n-1})}

für $x = r$ verschwindet. Da aber $r$ jede Wurzel von $X_n$ sein kann, und
diese Wurzeln alle von einander verschieden sind, so muss $\varPhi(x)$ durch
$X_n$ theilbar sein, also

\weq{4}{\varPhi(x) = X_n\,\varPsi(x).}

Die Function $\varPhi(x)$ hat, wie aus ihrer Bildungsweise hervorgeht, den
Co\"efficienten $1$ der höchsten Potenz von $x$ und sonst ganzzahlige
Co\"efficienten; ebenso $X_n$. Es muss also auch $\varPsi(x)$ den
Co\"efficienten $1$ der höchsten Potenz von $x$ und sonst ganzzahlige
Co\"efficienten haben. Es folgt aber aus (2) und (3), dass $\varPhi(x)$ für
$x = 1$ den Werth $\pm 1$ erhält, also

\weq{5}{\pm 1 = X_n(1)\,\varPsi(1),}

worin $\varPsi(1)$ eine ganze Zahl sein muss.

Wenn nun $n$ eine Primzahlpotenz $p^{\pi}$ ist, so ist $X_n(1) = p$, und dies
widerspricht der Gleichung (5). Also ist die Irreducibilität von $X_n$ für
den Fall bewiesen, dass $n$ eine Primzahl oder eine Potenz einer Primzahl
ist.

Der allgemeine Beweis verlangt also noch Folgendes: \gesperrt{Sind $a$ und
$b$ zwei ganze Zahlen ohne gemeinsamen Theiler, wird ferner vorausgesetzt,
dass $X_a$ und $X_b$ irreducibel sind, so folgt, dass auch $X_{ab}$
irreducibel ist.}

Um dies zu beweisen, setzen wir $n = ab$ und nehmen an, es sei $\varphi(x)$
ein rationaler Factor von $X_n$. Die Wurzeln von $\varphi(x) = 0$ sind dann
unter den primitiven $n^{\mathrm{ten}}$ Einheitswurzeln enthalten, und sind
also von der Form

\weq{6}{\varrho = \alpha\beta,}

worin $\alpha$ und $\beta$ gewisse primitive $a^{\mathrm{te}}$ und
$b^{\mathrm{te}}$ Einheitswurzeln bedeuten (\S.\,132).

Nun lässt sich (durch einen speciellen Fall der Tschirnhausen-Transformation)
aus $\varphi(x)$ eine Function $\varPhi(x)$ mit rationalen Co\"efficienten
ableiten, deren Wurzeln die $a^{\mathrm{ten}}$ Potenzen der Wurzeln von
$\varphi(x)$ sind. Es sind also die Wurzeln von $\varPhi(x)$
\[
\varrho^a = \beta^a,
\]
d. h., da $a$ relativ prim zu $b$ ist, lauter primitive $b^{\mathrm{te}}$
Einheitswurzeln, und wenn also überhaupt ein $\varrho$ vorhanden ist, so muss
$\varPhi(x)$ mit $X_b$ einen nicht constanten Factor gemein haben. Nun haben
wir aber als schon bewiesen vorausgesetzt, dass $X_b$ keinen rationalen
Factor ausser sich selbst habe; also muss $\varPhi(x)$ durch
%CONT>>
